% \iffalse meta-comment
%
% File: ltx-talk-overlay.dtx Copyright (C) 2024-2026 Joseph Wright
%
% It may be distributed and/or modified under the conditions of the
% LaTeX Project Public License (LPPL), either version 1.3c of this
% license or (at your option) any later version.  The latest version
% of this license is in the file
%
%    https://www.latex-project.org/lppl.txt
%
% This file is part of the "ltx-talk bundle" (The Work in LPPL)
% and all files in that bundle must be distributed together.
%
% The released version of this bundle is available from CTAN.
%
% -----------------------------------------------------------------------
%
% The development version of the bundle can be found at
%
%    https://github.com/josephwright/ltx-talk
%
% for those people who are interested.
%
% -----------------------------------------------------------------------
%
%<*driver>
\documentclass{l3doc}
% Additional commands needed in this source
\NewDocumentCommand\email{m}{\href{mailto:#1}{\nolinkurl{#1}}}
\begin{document}
  \DocInput{\jobname.dtx}
\end{document}
%</driver>
% \fi
%
% ^^A As we are dealing with a class, this has to be done manually
% \def\filedate{2026-11-05}
% \def\fileversion{v0.6.11}
%
% \title{^^A
%   \pkg{ltx-talk-overlay} -- Overlays^^A
%   \thanks{This file describes \fileversion,
%     last revised \filedate.}^^A
% }
%
% \author{^^A
%  Joseph Wright^^A
%  \thanks{^^A
%    E-mail: \email{joseph@texdev.net}^^A
%   }^^A
% }
%
% \date{Released \filedate}
%
% \maketitle
%
% \begin{documentation}
%
% \end{documentation}
%
% \begin{implementation}
%
% \section{\pkg{ltx-talk-overlay} implementation}
%
% Start the \pkg{DocStrip} guards.
%    \begin{macrocode}
%<*class>
%    \end{macrocode}
%
% Identify the internal prefix.
%    \begin{macrocode}
%<@@=talk>
%    \end{macrocode}
%
% \subsection{Utilities}
%
% \begin{macro}[TF]{\@@_if_overlay:n, \@@_if_overlay:V}
% \begin{macro}{\@@_overlay_arg:n}
%    \begin{macrocode}
\prg_new_protected_conditional:Npnn \@@_if_overlay:n #1 { T , F , TF }
  {
    \@@_decode_parse:n {#1}
    \bool_if:NTF \l_@@_decode_overlays_bool
      \prg_return_true:
      \prg_return_false:
  }
\prg_generate_conditional_variant:Nnn \@@_if_overlay:n { V } { T , F , TF }
%    \end{macrocode}
%   A macro processor variant of the check that always results in an |N|-type
%   bool.
%    \begin{macrocode}
\cs_new_protected:Npn \@@_overlay_arg:n #1
  {
    \@@_if_overlay:nTF {#1}
      { \cs_set:Npn \ProcessedArgument { \c_true_bool } }
      { \cs_set:Npn \ProcessedArgument { \c_false_bool } }
  }
%    \end{macrocode}
% \end{macro}
% \end{macro}
%
% \begin{variable}{\l_@@_shuffle_skip}
%   For tracking.
%    \begin{macrocode}
\skip_new:N \l_@@_shuffle_skip
%    \end{macrocode}
% \end{variable}
%
% \begin{macro}{\@@_shuffle_skip:n}
%   As opacity uses whatsits at present, we need to make sure that any spaces
%   come \emph{after} them. This is done by \enquote{shuffling} the last skip
%   past the opacity.
%    \begin{macrocode}
\cs_new_protected:Npn \@@_shuffle_skip:n #1
  {
    \skip_set_eq:NN \l_@@_shuffle_skip \tex_lastskip:D
    \bool_lazy_and:nnTF
      { ! \skip_if_eq_p:nn \l_@@_shuffle_skip { 0pt } }
      {
        \bool_lazy_or_p:nn
          { \mode_if_horizontal_p: }
          { \mode_if_vertical_p: }
      }
      {
        \tex_unskip:D
        #1
        \mode_if_horizontal:TF
          { \skip_horizontal:n }
          { \skip_vertical:n }
            \l_@@_shuffle_skip
      }
      {#1}
  }
%    \end{macrocode}
% \end{macro}
%
% \subsection{Opacity utilities}
%
% Currently, opacity is applies using whatsits at a low level. That means
% that to preserve spacing, we need to insert no-op versions in various
% places. To do that and get correct overlays, we need to track the current
% opacity. At present, this seems very \cls{ltx-talk}-specific, so
% is handled here with a few auxiliaries.
%
% \begin{variable}{\g_@@_opacity_seq, \g_@@_hidden_bool}
%   The opacities in force, innermost first, and whether the innermost is
%   zero, which is what makes material hidden rather than dimmed (see the
%   section on hidden text below).
%    \begin{macrocode}
\seq_new:N \g_@@_opacity_seq
\bool_new:N \g_@@_hidden_bool
%    \end{macrocode}
% \end{variable}
%
% \begin{macro}{\@@_hidden_update:}
%   Sets the flag from the innermost opacity. It is global: an opacity set
%   by \cs{pause} lasts until the end of the slide, past the end of any
%   group it was set in. The values are stored as decimals, so a string
%   comparison is enough, allowing for the sign that a zero can keep.
%    \begin{macrocode}
\cs_new_protected:Npn \@@_hidden_update:
  {
    \seq_get:NNTF \g_@@_opacity_seq \l_@@_tmp_tl
      {
        \bool_gset:Nn \g_@@_hidden_bool
          {
               \str_if_eq_p:Vn \l_@@_tmp_tl { 0 }
            || \str_if_eq_p:Vn \l_@@_tmp_tl { -0 }
          }
      }
      { \bool_gset_false:N \g_@@_hidden_bool }
  }
%    \end{macrocode}
% \end{macro}
%
% \begin{macro}{\@@_opacity_begin:n, \@@_opacity_end:}
%   Tracking wrappers.
%    \begin{macrocode}
\cs_new_protected:Npn \@@_opacity_begin:n #1
  {
    \seq_gpush:Ne \g_@@_opacity_seq { \fp_to_decimal:n {#1} }
    \@@_hidden_update:
    \@@_shuffle_skip:n { \opacity_begin:n {#1} }
  }
\cs_new_protected:Npn \@@_opacity_end:
  {
    \seq_gpop:NN \g_@@_opacity_seq \l_@@_tmp_tl
    \@@_hidden_update:
    \@@_shuffle_skip:n { \opacity_end: }
  }
%    \end{macrocode}
% \end{macro}
%
% \subsection{Action commands and environments}
%
% Commands that can be used as actions all have a common form (with one
% exception). The common internal structure is used to enable them to be used
% as actions by looking for the name \cs[no-index]{@@_action_\meta{name}:N}.
%
% \begin{macro}{\@@_action_alert:N}
%   At present a color selection.
%    \begin{macrocode}
\cs_new_protected:Npn \@@_action_alert:N #1
  {
    \bool_if:NTF #1
      { \color_select:n { alert } }
      { \color_select:n { . } }
  }
%    \end{macrocode}
% \end{macro}
%
% \begin{macro}
%   {
%     \@@_action_invisible:N     ,
%     \@@_action_invisible_end:N ,
%     \@@_action_visible:N       ,
%     \@@_action_visible_end:N
%   }
%   Simply (un)hide unconditionally, overwriting any previous opacity.
%    \begin{macrocode}
\cs_new_protected:Npn \@@_action_invisible:N #1
  {
    \bool_if:NTF #1
      { \@@_opacity_begin:n { 0 } }
      { \@@_opacity_begin:n { 1 } }
  }
\cs_new_protected:Npn \@@_action_invisible_end:N #1
  { \@@_opacity_end: }
\cs_new_protected:Npn \@@_action_visible:N #1
  {
    \bool_if:NTF #1
      { \@@_opacity_begin:n { 1 } }
      { \@@_opacity_begin:n { 0 } }
  }
\cs_new_protected:Npn \@@_action_visible_end:N #1
  { \@@_opacity_end: }
%    \end{macrocode}
% \end{macro}
%
% \begin{macro}{\@@_action_only:N, \@@_action_only_end:N}
% \begin{macro}{\@@_saved_ifendpe:}
%   Here, we simply throw away the content we do not want: this is done by
%   typesetting in a disposable box. As box setting rather than insertion
%   creates tag structures, they may need to be suspended. To avoid issues
%   with whatsits, we have a surrounding horizontal box if required. Suspending
%   tagging does not turn off or save/restore the \texttt{@endpe} mechanism,
%   so we need to do this ourselves. This avoids issues with paragraph tagging
%   if the mechanism would otherwise be active.
%    \begin{macrocode}
\cs_new_protected:Npn \@@_action_only:N #1
  {
    \bool_if:NF #1
      {
        \bool_if:NT \g_@@_frame_tag_bool
          {
            \cs_set_eq:NN \@@_saved_ifendpe: \if@endpe
            \vbox_set:Nw \l_@@_tmp_box
              \hbox_set:Nw \l_@@_tmp_box
                \tag_suspend:n { onlyenv }
          }
        \vbox_set:Nw \l_@@_tmp_box
      }
  }
\cs_new_protected:Npn \@@_action_only_end:N #1
  {
    \bool_if:NF #1
      {
        \vbox_set_end:
        \bool_if:NT \g_@@_frame_tag_bool
          {
                \tag_resume:n { onlyenv }
              \hbox_set_end:
            \vbox_set_end:
            \cs_gset_eq:NN \if@endpe \@@_saved_ifendpe:
          }
      }
  }
\cs_new_eq:NN \@@_saved_ifendpe: \if@endpe
%    \end{macrocode}
% \end{macro}
% \end{macro}
%
% \begin{variable}{\l_@@_uncover_hidden_fp}
%   Currently just an on-off, but that will change.
%    \begin{macrocode}
\NewTemplateType { hidden } { 0 }
\DeclareTemplateInterface { hidden } { talk } { 0 }
  { opacity : real = 0 }
\DeclareTemplateCode { hidden } { talk } { 0 }
  { opacity = \l_@@_uncover_hidden_fp }
  { \@@_opacity_begin:n { \l_@@_uncover_hidden_fp } }
\DeclareInstance { hidden } { std } { talk } { }
%    \end{macrocode}
% \end{variable}
% \begin{macro}{\@@_action_uncover:N, \@@_action_uncover_end:N}
%   Use the template: we may need to extend that to deal with the
%   end-of-template case later.
%    \begin{macrocode}
\cs_new_protected:Npn \@@_action_uncover:N #1
  {
    \bool_if:NTF #1
      { \@@_opacity_begin:n { 1 } }
      { \UseInstance { hidden } { std } }
  }
\cs_new_protected:Npn \@@_action_uncover_end:N #1
  { \@@_opacity_end: }
%    \end{macrocode}
% \end{macro}
%
% \begin{macro}{\invisible, \uncover, \visible}
%   All generated automatically using the above implementations.
%    \begin{macrocode}
\clist_map_inline:nn { invisible , uncover , visible }
  {
    \ExpandArgs { cne } \NewDocumentCommand {#1}
      { > { \@@_overlay_arg:n } D <> { all } +m }
      {
        \exp_not:c { @@_action_ #1 :N } ##1
        ##2
        \exp_not:c { @@_action_ #1 _end:N } ##1
      }
%    \end{macrocode}
% \end{macro}
% \begin{environment}{invisibleenv, uncoverenv, visibleenv}
%   And the environment versions.
%    \begin{macrocode}
    \ExpandArgs { nnee } \NewDocumentEnvironment { #1 env }
      { > { \@@_overlay_arg:n } D <> { all } }
      { \exp_not:c { @@_action_ #1 :N } ##1 }
      { \exp_not:c { @@_action_ #1 _end:N } ##1 }
  }
%    \end{macrocode}
% \end{environment}
%
% \begin{macro}{\alert}
%   The \cs{alert} command requires a group to contain color, so is done
%   separately even though it still uses basically the same mechanism.
%    \begin{macrocode}
\NewDocumentCommand \alert { > { \@@_overlay_arg:n } D <> { all } +m }
  {
    \group_begin:
      \@@_action_alert:N #1
        #2
     \group_end:
  }
%    \end{macrocode}
% \end{macro}
% \begin{environment}{alertenv}
%   As does the environment.
%    \begin{macrocode}
\NewDocumentEnvironment { alertenv } { > { \@@_overlay_arg:n } D <> { all } }
  { \@@_action_alert:N #1 }
  { }
%    \end{macrocode}
% \end{environment}
% \begin{macro}{\only}
%   This code needs to be done manually as for the command version the content
%   must be entirely discarded. That can't work for the environment version,
%   which has to deal with for example single items in a list (and so cannot
%   be collected up verbatim and must use a box).
%    \begin{macrocode}
\NewDocumentCommand \only { D <> { all } +m }
  {
    \@@_if_overlay:nT {#1}
      {#2}
  }
%    \end{macrocode}
% \end{macro}
% \begin{environment}{onlyenv}
%   The environment version could be done above, but it is clearer to keep
%   this code entirely separate from the rest.
%    \begin{macrocode}
\NewDocumentEnvironment { onlyenv } { > { \@@_overlay_arg:n } D <> { all } }
  { \@@_action_only:N #1 }
  { \@@_action_only_end:N #1 }
%    \end{macrocode}
% \end{environment}
%
% \begin{variable}
%   {
%     \l_@@_saved_overlays_bool ,
%     \l_@@_saved_action_str    ,
%     \l_@@_saved_actions_bool
%   }
%    \begin{macrocode}
\bool_new:N \l_@@_saved_overlays_bool
\str_new:N \l_@@_saved_action_str
\bool_new:N \l_@@_saved_actions_bool
%    \end{macrocode}
% \end{variable}
%
% \begin{variable}
%   {\l_@@_overlay_all_bool}
%    \begin{macrocode}
\bool_new:N \l_@@_overlay_all_bool
%    \end{macrocode}
% \end{variable}
% \begin{macro}{\action}
% \begin{environment}{actionenv}
% \begin{macro}{\@@_action_begin:n}
% \begin{macro}{\@@_action_begin:w}
% \begin{macro}{\@@_action_begin_auxi:n, \@@_action_begin_auxii:n}
% \begin{macro}{\@@_action_end:}
%   As we need data on not just overlays but also actions at the end of the
%   environment, this has to be done manually. To allow working with
%   environments but also items, the code needs to save data for the end
%   function. The group is needed for cases where we are not in a \LaTeX{}
%   environment group. When an \cs{onslide}/\cs{pause} is active, it takes
%   priority: sorted by applying up-front. Actions can be skipped entirely
%   if the overlay spec is simply \texttt{all}, as there will never be
%   any spacing issues, \foreign{etc}.
%    \begin{macrocode}
\NewDocumentCommand \action { d <> +m }
  {
    \group_begin:
      \@@_action_begin:n {#1}
      #2
      \@@_action_end:
    \group_end:
  }
\NewDocumentEnvironment { actionenv } { d <> }
  { \@@_action_begin:n {#1} }
  { \@@_action_end: }
\cs_new_protected:Npn \@@_action_begin:n #1
  {
    \group_begin:
      \tl_if_novalue:nTF {#1}
        {
          \exp_after:wN \@@_action_begin:w
            \l_@@_action_spec_str < all > \q_stop
        }
        { \@@_action_begin_auxi:n {#1} }
  }
\cs_new_protected:Npn \@@_action_begin:w #1 < #2 > #3 \q_stop
  { \@@_action_begin_auxi:n {#2} }
\cs_new_protected:Npn \@@_action_begin_auxi:n #1
  {
    \str_if_eq:nnTF {#1} { all }
      { \bool_set_true:N \l_@@_overlay_all_bool }
      {
        \bool_set_false:N \l_@@_overlay_all_bool
        \@@_action_begin_auxii:n {#1}
      }
  }
\cs_new_protected:Npn \@@_action_begin_auxii:n #1
  {
      \@@_decode_parse:n {#1}
      \bool_set_eq:NN \l_@@_saved_overlays_bool
        \l_@@_decode_overlays_bool
      \str_set_eq:NN \l_@@_saved_action_str
        \l_@@_decode_action_str
      \bool_set_eq:NN \l_@@_saved_actions_bool
        \l_@@_decode_actions_bool
      \tl_if_empty:NTF \g_@@_onslide_tl
        {
          \bool_if:NTF \l_@@_decode_overlays_bool
            {
              \cs_if_exist_use:cF
                { @@_action_ \l_@@_decode_action_str :N }
                  { \use_none:n }
                    \l_@@_decode_actions_bool
            }
            { \UseInstance { hidden } { std } }
        }
        { \@@_action_invisible:N \c_true_bool }
  }
\cs_new_protected:Npn \@@_action_end:
  {
      \bool_if:NF \l_@@_overlay_all_bool
        {
          \tl_if_empty:NTF \g_@@_onslide_tl
            {
              \bool_if:NTF \l_@@_saved_overlays_bool
                {
                  \cs_if_exist_use:cF
                    { @@_action_ \l_@@_saved_action_str _end:N }
                    { \use_none:n }
                      \l_@@_saved_actions_bool
                }
                { \@@_opacity_end: }
            }
            { \@@_action_invisible_end:N \c_true_bool }
        }
    \group_end:
  }
%    \end{macrocode}
% \end{macro}
% \end{macro}
% \end{macro}
% \end{macro}
% \end{environment}
% \end{macro}
%
% \subsection{Non-action commands and environments}
%
% This section contains commands and environments that do \emph{not}
% need to be made available as actions.
%
% \begin{macro}{\alt}
%   Simple wrappers around the internal switch.
%    \begin{macrocode}
\NewDocumentCommand \alt { D <> { all } +m +m }
  {
    \@@_if_overlay:nTF {#1}
      {#2}
      {#3}
  }
%    \end{macrocode}
% \end{macro}
%
% \begin{macro}{\onslide}
% \begin{macro}{\@@_onslide:n}
%   Simply make transparent: this is done without grouping so we can work
%   for example in tabular cells.
%    \begin{macrocode}
\NewDocumentCommand \onslide { D <> { all } }
  {
    \@@_onslide:n {#1}
    \ignorespaces
  }
\cs_new_protected:Npn \@@_onslide:n #1
  {
    \tl_use:N \g_@@_onslide_tl
    \tl_gclear:N \g_@@_onslide_tl
    \@@_if_overlay:nF {#1}
      {
        \@@_opacity_begin:n { 0 }
        \tl_gput_right:Nn \g_@@_onslide_tl
          { \@@_opacity_end: }
      }
  }
%    \end{macrocode}
% \end{macro}
% \end{macro}
% \begin{variable}{\g_@@_onslide_tl}
%    \begin{macrocode}
\tl_new:N \g_@@_onslide_tl
%    \end{macrocode}
% \end{variable}
%
% \begin{macro}{\temporal}
%   A tricky one: to separate the not-on-current-slide cases, the flag to
%   continue is used.
%    \begin{macrocode}
\NewDocumentCommand \temporal { D <> { all } +m +m +m }
  {
    \@@_if_overlay:nTF {#1}
      {#3}
      {
        \bool_if:NTF \l_@@_decode_trailing_bool
          {#4}
          {#2}
      }
  }
%    \end{macrocode}
% \end{macro}
%
% \begin{macro}{\pause}
%   A thin wrapper.
%    \begin{macrocode}
\NewDocumentCommand \pause { o }
  {
    \legacy_if:nF { measuring@ }
      {
        \IfNoValueTF {#1}
          { \int_gincr:N \g_@@_pauses_int }
          { \int_gset:Nn \g_@@_pauses_int {#1} }
        \exp_args:Ne \@@_onslide:n
          { \int_eval:n { \g_@@_pauses_int + 1 } - }
      }
  }
%    \end{macrocode}
% \end{macro}
%
% \subsection{Fixed-size areas}
%
% \begin{macro}{\@@_overprint_begin:n}
%   A common auxiliary for overprinting, which starts off much the same for
%   both \env{overlayarea} and \env{overprint}.
%    \begin{macrocode}
\cs_new_protected:Npn \@@_overprint_begin:n #1
  {
    \par
    \vbox_set_to_wd:Nnw \l_@@_tmp_box {#1}
      \raggedright
      \ignorespaces
  }
%    \end{macrocode}
% \end{macro}
%
% \begin{environment}{overlayarea}
%   An initial approach: quite similar to a column.
%    \begin{macrocode}
\NewDocumentEnvironment { overlayarea } { m m }
  { \@@_overprint_begin:n {#1} }
  {
    \vbox_set_end:
    \vbox_to_ht:nn {#2}
      {
        \box_use_drop:N \l_@@_tmp_box
        \vfil
      }
    \par
  }
%    \end{macrocode}
% \end{environment}
%
% \begin{variable}{\l_@@_overprint_int}
%   Track the overprints on a slide: as the slide forms a group, we do not need
%   to worry about resetting.
%    \begin{macrocode}
\int_new:N \l_@@_overprint_int
%    \end{macrocode}
% \end{variable}
%
% \begin{macro}[EXP]{\@@_frame_overprint:}
%   To refer to the current overprint environment within the document:
%   needed in the \texttt{.aux} so avoids using non-letters.
%    \begin{macrocode}
\cs_new:Npn \@@_frame_overprint:
  {
    \int_to_Roman:n \g_@@_frame_int
    \int_to_roman:n \l_@@_overprint_int
  }
%    \end{macrocode}
% \end{macro}
%
% \begin{environment}{overprint}
% \begin{macro}{\@@_overprint_save_ht:}
% \begin{macro}{\@@_overprint_check_ht:n}
%   For overprinting, in contrast to \cls{beamer} we use a two-pass approach
%   to save the size at the end of the run: this means you can use \cs{only}
%   for example in overprinting.
%    \begin{macrocode}
\NewDocumentEnvironment { overprint } { O { \textwidth } }
  { \@@_overprint_begin:n {#1} }
  {
    \vbox_set_end:
    \int_incr:N \l_@@_overprint_int
    \@@_overprint_save_ht:
    \cs_if_exist:cTF 
      { overprint@ \@@_frame_overprint: }
      {
        \dim_compare:vNnTF
          { overprint@ \@@_frame_overprint: }
            > { \box_ht:N \l_@@_tmp_box }
          {
            \vbox_to_ht:vn
              { overprint@ \@@_frame_overprint: }
              {
                \box_use_drop:N \l_@@_tmp_box
                \vfil
              }
          }
          { \box_use_drop:N \l_@@_tmp_box }
      }
      { \box_use_drop:N \l_@@_tmp_box }
    \par
  }
%    \end{macrocode}
%   As there is no clear end-point for overprinting, we need to be careful to
%   keep the current width separate from the saved one. The rest is then about
%   saving to the \texttt{.aux} file and helping out the user. 
%    \begin{macrocode}
\cs_new_protected:Npn \@@_overprint_save_ht:
  {
    \tl_if_exist:cF { g_@@_overprint_ \@@_frame_overprint: _tl }
      {
        \tl_new:c { g_@@_overprint_ \@@_frame_overprint: _tl }
        \tl_gset:cn { g_@@_overprint_ \@@_frame_overprint: _tl }
          { 0pt }
      }
    \tl_gset:ce { g_@@_overprint_ \@@_frame_overprint: _tl }
      {
        \dim_max:vn { g_@@_overprint_ \@@_frame_overprint: _tl }
          { \box_ht:N \l_@@_tmp_box }
      }
    \legacy_if:nT { @filesw }
      {
        \iow_now:Ne \@auxout
          {
            \gdef \exp_not:c { overprint@ \@@_frame_overprint: }
              {
                \exp_not:v { g_@@_overprint_ \@@_frame_overprint: _tl }
              }
          }
      }
    \hook_gput_code:nne { enddocument / afterlastpage } { talk }
       { \@@_overprint_check_ht:n { \@@_frame_overprint: } }
  }
\cs_new_protected:Npn \@@_overprint_check_ht:n #1
  {
    \bool_lazy_and:nnF
      { \exp_not:N \cs_if_exist_p:c { overprint@ #1 } }
      {
        \dim_compare_p:vNv { overprint@ #1 } = { g_@@_overprint_ #1 _tl }
      }
      {
        \msg_warning:nn { talk } { overprint-ht }
        \cs_gset_protected:Npn \@@_overprint_check_ht:n ##1 { }
      }
  }
\msg_new:nnn { talk } { overprint-ht }
  {
    Overprint~area~height~has~changed:\\
    rerun~LaTeX.
  }
%    \end{macrocode}
% \end{macro}
% \end{macro}
% \end{environment}
%
% \subsection{Adding overlays to existing commands}
%
% \begin{macro}
%  {
%    \textbf ,
%    \textit ,
%    \textmd ,
%    \textnormal ,
%    \textrm ,
%    \textsc ,
%    \textsf ,
%    \textsl ,
%    \texttt ,
%    \textup ,
%    \emph
%  }
% \begin{macro}
%  {
%    \stdtextbf ,
%    \stdtextit ,
%    \stdtextmd ,
%    \stdtextnormal ,
%    \stdtextrm ,
%    \stdtextsc ,
%    \stdtextsf ,
%    \stdtextsl ,
%    \stdtexttt ,
%    \stdtextup ,
%    \stdemph
%  }

% \begin{macro}{\@@_textcmd_eqiv:n}
%   Make the standard text commands overlay-aware. To keep the spacing
%   unchanged when the command is not active, we use the same approach as
%   the kernel does for inserting the right grouping.
%    \begin{macrocode}
\tl_map_inline:nn
  {
    \textbf
    \textit
    \textmd
    \textnormal
    \textrm
    \textsc
    \textsf
    \textsl
    \texttt
    \textup
    \emph
  }
  {
    \ExpandArgs { c } \NewCommandCopy { std \cs_to_str:N #1 } #1
    \ExpandArgs { Nne } \RenewDocumentCommand #1
      { D <> { all } +m }
      {
        \exp_not:N \@@_if_overlay:nTF {##1}
          { \exp_not:c { std \cs_to_str:N #1 } }
          { \exp_not:N \@@_textcmd_eqiv:n }
            {##2}
      }
  }
\cs_new_protected:Npn \@@_textcmd_eqiv:n #1
  {
    \mode_if_math:TF
      { { \mbox {#1} } }
      {
        \mode_leave_vertical:
        {#1}
      }
  }
%    \end{macrocode}
% \end{macro}
% \end{macro}
% \end{macro}
%
% \begin{macro}{\includegraphics}
% \begin{macro}{\stdincludegraphics}
%   Just wrap up the args and forward if appropriate. The star is |#1| here as
%   that matches the documented behavior of starred commands generally.
%    \begin{macrocode}
\RequirePackage { graphicx }
\NewCommandCopy \stdincludegraphics \includegraphics
\RenewDocumentCommand \includegraphics { s D <> { all } o o m }
  {
    \@@_if_overlay:nT {#2}
      {
        \use:e
          {
            \exp_not:N \stdincludegraphics
              \IfBooleanT #1 { * }
              \IfNoValueF {#3} { [ \exp_not:n { {#3} } ] }
              \IfNoValueF {#4} { [ \exp_not:n { {#4} } ] }
          }
            {#5}
      }
  }
%    \end{macrocode}
% \end{macro}
% \end{macro}
%
% \begin{macro}{\label}
% \begin{macro}{\@@_label:n}
%   Here, we can't wrap the existing command up as we need the space hack, so
%   it has to be declared from scratch. There is also a non-standard overlay
%   default. At present, no special tricks as seen in \pkg{beamer}.   
%    \begin{macrocode}
\RenewDocumentCommand \label { D <> { 1 } m }
  {
    \@bsphack
    \@@_if_overlay:nT {#1}
      { \@@_label:n {#2} }
    \@esphack
  }
\cs_new_protected:Npn \@@_label:n #1
  {
    \begingroup
      \UseHookWithArguments { label } { 1 } {#1}
      \protected@write \@auxout { }
        {
          \string \newlabel {#1}
            { 
              { \@currentlabel }
              { \thepage }
              { \@currentlabelname }
              { \@currentHref }
              { \@kernel@reserved@label@data }
            }
        }
    \endgroup
  }
%    \end{macrocode}
% \end{macro}
% \end{macro}
% \begin{macro}{\std@label@in@display, \std@label@in@display}
%   We also need to cover \cs{label@in@display}: a bit more awkward as this is
%   a document command but not set up as such in \pkg{amsmath}. For the present,
%   cross our fingers on this!
%    \begin{macrocode}
\cs_new_eq:NN \std@label@in@display \label@in@display
\RenewDocumentCommand \label@in@display { D <> { 1 } m }
  {
    \@@_if_overlay:nT {#1}
      { \std@label@in@display {#2} }
  }
%    \end{macrocode}
% \end{macro}
%
% \subsection{Overlay patching of third-party environments}
%
% To allow opacity to apply to the entirety of constructed graphics, we need to
% apply a transparency group around the whole thing. We need to do that by saving
% the input in an Xform object, which then allows the group to be applied.
%
% \begin{variable}{\g_@@_opacity_group_int}
%   Each use needs an Xform object, which at present we have to track as
%   there are no anonymous ones.
%    \begin{macrocode}
\int_new:N \g_@@_opacity_group_int
%    \end{macrocode}
% \end{variable}
%
% \begin{macro}{\@@_opacity_group_begin:, \@@_opacity_group_end:}
%   A general mechanism to collect up an environment in a box, which we can
%   then use as the source for an Xform object.
%    \begin{macrocode}
\cs_new_protected:Npn \@@_opacity_group_begin:
  { \begin { lrbox } \l_@@_tmp_box }
\cs_new_protected:Npn \@@_opacity_group_end:
  {
    \end { lrbox }
    \mode_leave_vertical:
    \int_gincr:N \g_@@_opacity_group_int
    \exp_args:Ne \pdfxform_new:nnn
      { talk.opacity \int_use:N \g_@@_opacity_group_int }
      { /Group << /S /Transparency /K ~ false /I ~ false >> }
      { \usebox \l_@@_tmp_box }
    \exp_args:Ne \pdfxform_use:n
      { talk.opacity \int_use:N \g_@@_opacity_group_int }
  }
%    \end{macrocode}
% \end{macro}
%
% At present, we only add code onto the main top-level functions. This is only
% applied in Lua\TeX{} due to restrictions on Xform structures in \pdfTeX{}.
% Here, we apply the collection code to the commands not the environment
% forms to deal with \enquote{direct} usage.
%    \begin{macrocode}
\sys_if_engine_luatex:T
  {
    \AddToHook { cmd / pspicture / before } { \@@_opacity_group_begin: }
    \AddToHook { cmd / endpspicture / after } { \@@_opacity_group_end: }
  }
%    \end{macrocode}
%
% \subsection{Hidden text}
%
% Content hidden by an overlay is drawn at opacity~$0$, which leaves it in
% the page: a PDF reader's search and text selection find its text, and on
% a tagged slide its structure is there for assistive technology to read.
% Two things are done about that, both selected by the
% |remove-hidden-text| option.
%
% A formula is wholly inside or wholly outside a hidden region, so a
% formula that starts while hidden is made an artifact as a whole: it gets
% no structure element, alternative text or MathML. That is done through
% the math tagging sockets, so it works with any engine; the flag set by
% \cs{@@_hidden_update:} says whether a formula starts hidden. Hidden plain
% text keeps its structure, but with Lua\TeX{} its glyphs are removed from
% the page just before shipout: the code in \file{ltx-talk.lua} replaces
% every glyph drawn at opacity~$0$ by a kern of the same width, so the
% layout is unchanged and the text is gone. Rules, images and drawings are
% left alone, as they carry no text.
%
% For the glyphs, what counts is where they are drawn, not when they were
% typeset, and so the walk follows the opacity operators on the page
% rather than a flag or attribute set during typesetting. The two differ
% for material built before the opacity changes but placed after it: the
% label of an \cs{item} with an overlay specification is one, as it is
% boxed before the item's action switches the opacity, and it would keep
% its text if only nodes typeset while hidden were removed.
%
% \begin{variable}{\g_@@_formula_seq}
%   One entry per formula being typeset, recording whether it was made an
%   artifact, so that its end matches its start even if the opacity has
%   changed in between.
%    \begin{macrocode}
\seq_new:N \g_@@_formula_seq
%    \end{macrocode}
% \end{variable}
%
% \begin{macro}{\@@_formula_hide:, \@@_formula_hide_set:}
%   A hidden formula switches paragraph tagging off, as the standard plug
%   does, and \pkg{luamml}, when loaded, is told to ignore it. Which of the
%   two commands exist is settled once, when the plugs are installed, by
%   expansion.
%    \begin{macrocode}
\cs_new_protected:Npn \@@_formula_hide: { }
\cs_new_protected:Npn \@@_formula_hide_set:
  {
    \cs_gset_protected:Npe \@@_formula_hide:
      {
        \cs_if_exist:NT \tagpdfparaOff
          { \exp_not:N \tagpdfparaOff }
        \cs_if_exist:NT \luamml_ignore:
          { \exp_not:N \luamml_ignore: }
      }
  }
%    \end{macrocode}
% \end{macro}
%
% \begin{macro}{\@@_formula_begin:nn, \@@_formula_end:n}
%   The begin auxiliary takes the code of the plug that was in place and
%   the formula; the end auxiliary takes the code of the plug that was in
%   place. A hidden formula is wrapped in an artifact marked-content
%   sequence.
%    \begin{macrocode}
\cs_new_protected:Npn \@@_formula_begin:nn #1#2
  {
    \bool_if:NTF \g_@@_hidden_bool
      {
        \seq_gpush:Nn \g_@@_formula_seq { hidden }
        \@@_formula_hide:
        \tag_mc_begin:n { artifact }
        #2
      }
      {
        \seq_gpush:Nn \g_@@_formula_seq { shown }
        #1
      }
  }
\cs_new_protected:Npn \@@_formula_end:n #1
  {
    \seq_gpop:NNF \g_@@_formula_seq \l_@@_tmp_tl
      { \tl_set:Nn \l_@@_tmp_tl { shown } }
    \str_if_eq:VnTF \l_@@_tmp_tl { hidden }
      { \tag_mc_end: }
      {#1}
  }
%    \end{macrocode}
% \end{macro}
%
% \begin{macro}{\@@_formula_plug_restore:nn}
%   Assigns to a formula socket the plug it had before the class's own,
%   which \cs{@@_formula_plugs:n} saved. The arguments are the socket name
%   after |math/| and the name of the saved plug's variable.
%    \begin{macrocode}
\cs_new_protected:Npn \@@_formula_plug_restore:nn #1#2
  {
    \exp_args:Nnv \AssignTaggingSocketPlug
      { math / #1 } { g_@@_math_ #2 _plug_tl }
  }
%    \end{macrocode}
% \end{macro}
%
% \begin{macro}{\@@_formula_plugs:n}
%   For |inline| or |display| math, installs plugs for the formula begin
%   and end sockets that hand over to the plugs in place before. That is
%   done at the end of the preamble, once the math tagging code has chosen
%   its plugs. To hand over, a plug assigns the previous plug to its socket,
%   uses the socket, and assigns itself again: the assignments are local,
%   and so they leave the socket as it was.
%    \begin{macrocode}
\cs_new_protected:Npn \@@_formula_plugs:n #1
  {
    \socket_get_plug:nN { tagsupport / math / #1 / formula / begin }
      \l_@@_tmp_tl
    \tl_gset_eq:cN { g_@@_math_ #1 _begin_plug_tl } \l_@@_tmp_tl
    \socket_get_plug:nN { tagsupport / math / #1 / formula / end }
      \l_@@_tmp_tl
    \tl_gset_eq:cN { g_@@_math_ #1 _end_plug_tl } \l_@@_tmp_tl
    \NewTaggingSocketPlug { math / #1 / formula / begin } { talk-hidden }
      {
        \@@_formula_begin:nn
          {
            \@@_formula_plug_restore:nn
              { #1 / formula / begin } { #1 _begin }
            \tag_socket_use:nnn { math / #1 / formula / begin }
              {##1} {##2}
            \AssignTaggingSocketPlug
              { math / #1 / formula / begin } { talk-hidden }
          }
          {##2}
      }
    \NewTaggingSocketPlug { math / #1 / formula / end } { talk-hidden }
      {
        \@@_formula_end:n
          {
            \@@_formula_plug_restore:nn
              { #1 / formula / end } { #1 _end }
            \tag_socket_use:n { math / #1 / formula / end }
            \AssignTaggingSocketPlug
              { math / #1 / formula / end } { talk-hidden }
          }
      }
    \AssignTaggingSocketPlug { math / #1 / formula / begin } { talk-hidden }
    \AssignTaggingSocketPlug { math / #1 / formula / end } { talk-hidden }
  }
%    \end{macrocode}
% \end{macro}
%
% With Lua\TeX{} the Lua code is always loaded. Both parts are opt-in:
% only with |remove-hidden-text| are the plugs installed and, with
% Lua\TeX{}, the removal of hidden glyphs switched on. Otherwise hidden
% content stays as it is typeset.
%    \begin{macrocode}
\sys_if_engine_luatex:T
  { \lua_load_module:n { ltx-talk } }
\bool_if:NT \l_@@_remove_hidden_text_bool
  {
    \hook_gput_code:nnn { begindocument / end } { talk }
      {
        \@@_formula_hide_set:
        \@@_formula_plugs:n { inline }
        \@@_formula_plugs:n { display }
      }
    \sys_if_engine_luatex:T
      { \lua_now:n { ltx.__talk.remove_hidden_text() } }
  }
%    \end{macrocode}
%
%    \begin{macrocode}
%</class>
%    \end{macrocode}
%
% \subsection{Lua code}
%
% With Lua\TeX{} the code here is always loaded, but it only acts once
% |ltx.__talk.remove_hidden_text()| has been called, which the class does for
% |remove-hidden-text|. The walk then runs in the
% |pre_shipout_filter| callback, on the finished page box.
%
%    \begin{macrocode}
%<*lua>
%    \end{macrocode}
%
% \begin{macro}{ltx.__talk}
%   The module's table, which is internal, following \pkg{tagpdf}'s
%   |ltx.__tag|.
%    \begin{macrocode}
luatexbase.provides_module({
  name        = "ltx-talk",
  date        = "2026-11-05",
  version     = "0.6.11",
  description = "Lua code for ltx-talk"
})
ltx = ltx or { }
ltx.__talk = ltx.__talk or { }
local talk = ltx.__talk
%    \end{macrocode}
% \end{macro}
%
%   Local copies of standard functions.
%    \begin{macrocode}
local insert_before = node.insert_before
local max           = math.max
local new_node      = node.new
local node_free     = node.free
local node_remove   = node.remove
local traverse      = node.traverse
%    \end{macrocode}
%
%   Node types, whatsit subtypes and the commands of a colour-stack
%   whatsit, all constants.
%    \begin{macrocode}
local id_disc            = node.id("disc")
local id_glue            = node.id("glue")
local id_glyph           = node.id("glyph")
local id_hlist           = node.id("hlist")
local id_kern            = node.id("kern")
local id_vlist           = node.id("vlist")
local id_whatsit         = node.id("whatsit")
local subtype_colorstack = node.subtype("pdf_colorstack")
local subtype_restore    = node.subtype("pdf_restore")
local subtype_save       = node.subtype("pdf_save")
local stack_set          = 0
local stack_push         = 1
local stack_pop          = 2
%    \end{macrocode}
%
% The opacity in force is tracked as a PDF reader would see it.
% \pkg{l3opacity} sets opacity through a colour stack in
% \enquote{page} mode, whose top value is reissued at the start of every
% page, and the operators it pushes name the value (|/opacity0 gs|). The
% stack is recognised by that data the first time it is pushed, and
% |opacity_stack| mirrors its contents across pages. Within a page,
% |opacity| is the value in force, and |saved_opacities| holds the values
% saved by |q| from the graphics drivers, to be restored by |Q|. The
% number of hidden glyphs met on each page is kept in
% |ltx.__talk.hidden_text_pages|, for tests.
%    \begin{macrocode}
local stack_number
local opacity_stack = { }
local opacity
local saved_opacities
local hidden_glyphs
talk.hidden_text_pages = { }
%    \end{macrocode}
%
% \begin{macro}[int]{opacity_value}
%   The opacity that the data of a colour-stack whatsit sets, or |nil|
%   for data that is not an opacity operator. Callers take |nil| to mean
%   fully visible.
%    \begin{macrocode}
local function opacity_value(data)
  local value = data and data:match("^%s*/opacity([%d%.]+)%s+gs%s*$")
  return value and tonumber(value)
end
%    \end{macrocode}
% \end{macro}
%
% \begin{macro}[int]{opacity_top}
%   The value at the top of the opacity stack, which is fully visible
%   when the stack is empty.
%    \begin{macrocode}
local function opacity_top()
  return opacity_stack[#opacity_stack] or 1
end
%    \end{macrocode}
% \end{macro}
%
% \begin{macro}[int]{track_opacity}
%   Updates the opacity in force for a whatsit: a push, pop or set on the
%   opacity stack, or a save or restore of the graphics state.
%    \begin{macrocode}
local function track_opacity(n)
  local subtype = n.subtype
  if subtype == subtype_colorstack then
    local value = opacity_value(n.data)
    if value and not stack_number then
      stack_number = n.stack
    end
    if n.stack == stack_number then
      local command = n.command
      if command == stack_push then
        opacity_stack[#opacity_stack + 1] = value or 1
      elseif command == stack_pop then
        opacity_stack[#opacity_stack] = nil
      elseif command == stack_set then
        opacity_stack[max(#opacity_stack,1)] = value or 1
      end
      opacity = opacity_top()
    end
  elseif subtype == subtype_save then
    saved_opacities[#saved_opacities + 1] = opacity
  elseif subtype == subtype_restore then
    if #saved_opacities > 0 then
      opacity = saved_opacities[#saved_opacities]
      saved_opacities[#saved_opacities] = nil
    end
  end
end
%    \end{macrocode}
% \end{macro}
%
% \begin{macro}[int]{replace_glyph}
%   Replaces a glyph in a list by a kern of the same width, allowing for
%   font expansion, and returns the new head of the list.
%    \begin{macrocode}
local function replace_glyph(head,glyph)
  local kern = new_node(id_kern,1)
  local width = glyph.width
  local expansion = glyph.expansion_factor
  if expansion and expansion ~= 0 then
    width = width + width * expansion // 1000000
  end
  kern.kern = width
  kern.attr = glyph.attr
  head = insert_before(head,glyph,kern)
  head = node_remove(head,glyph)
  node_free(glyph)
  return head
end
%    \end{macrocode}
% \end{macro}
%
% \begin{macro}[int]{remove_glyphs}
%   Walks a list, and the lists inside it, in the order they are drawn,
%   and replaces every glyph met at opacity~$0$. The glyphs are collected
%   first and replaced after the walk, so that the traversal is not
%   disturbed. After line breaking, the replace list of a discretionary is
%   what the page shows, but all three are walked. Rules and drawings met
%   at opacity~$0$ are left alone, as they carry no text.
%    \begin{macrocode}
local function remove_glyphs(head)
  local found
  for n in traverse(head) do
    local id = n.id
    if id == id_hlist or id == id_vlist then
      if n.list then
        n.list = remove_glyphs(n.list)
      end
    elseif id == id_glyph then
      if opacity == 0 then
        hidden_glyphs = hidden_glyphs + 1
        found = found or { }
        found[#found + 1] = n
      end
    elseif id == id_disc then
      if n.pre then
        n.pre = remove_glyphs(n.pre)
      end
      if n.post then
        n.post = remove_glyphs(n.post)
      end
      if n.replace then
        n.replace = remove_glyphs(n.replace)
      end
    elseif id == id_glue then
      local leader = n.leader
      if leader and (leader.id == id_hlist or leader.id == id_vlist)
        and leader.list then
        leader.list = remove_glyphs(leader.list)
      end
    elseif id == id_whatsit then
      track_opacity(n)
    end
  end
  if found then
    for _,glyph in ipairs(found) do
      head = replace_glyph(head,glyph)
    end
  end
  return head
end
%    \end{macrocode}
% \end{macro}
%
% \begin{macro}[int]{hidden_text_filter}
%   The callback function: starts each page at the opacity carried over
%   from the last one.
%    \begin{macrocode}
local function hidden_text_filter(box)
  opacity = opacity_top()
  saved_opacities = { }
  hidden_glyphs = 0
  if box.list then
    box.list = remove_glyphs(box.list)
  end
  talk.hidden_text_pages[status.total_pages + 1] = hidden_glyphs
  return true
end
%    \end{macrocode}
% \end{macro}
%
% \begin{macro}{ltx.__talk.remove_hidden_text}
%   Adds the callback. \pkg{tagpdf} marks the page's content in the same
%   callback, and the hidden glyphs must be gone first, or the marked
%   content sequences that held them are kept, holding only kerns.
%    \begin{macrocode}
function talk.remove_hidden_text()
  luatexbase.add_to_callback("pre_shipout_filter",hidden_text_filter,
    "ltx-talk.hidden-text")
  luatexbase.declare_callback_rule("pre_shipout_filter",
    "ltx-talk.hidden-text","before","tagpdf")
end
%    \end{macrocode}
% \end{macro}
%
%    \begin{macrocode}
%</lua>
%    \end{macrocode}
%
% \end{implementation}
%
% \PrintIndex
